\begin{theorem}\label{night:P3K.116}
\StatementRecord{night:P3K.116}{proved}{}{the original source, point line, roots, characters, reference, selected fan and scalar range stated here}
Retain the literal OS45 source, marked section and point line
\[
\begin{gathered}
 R=t^2+a,\qquad f=R^2+t,\qquad
 g=R^3+\frac32tR+\frac38,\\
 S_a:\ y^2=x^3-3fx+2g,\qquad
 G=(R,\sqrt3/2),\qquad M_8=\mathcal O([8G]-[O]).
\end{gathered}
\]
The suspension source map is translation by \(-G\); its simple zero is
the whole selected \(I_1\) fibre. Retain the original multiplicity-one
cap at the ordinary value \(t=0\), the complete prescribed line and
annular reference, and any prescribed selected lattice-basis fan.

The continuation of the initially prescribed directed line map is
fixed by its actual coefficient:
\[
 B_\lambda(p,v_M)=
 \bigl(p-G,\lambda C_*(t-t_c)B_{8,1}(p)v_M\bigr),\qquad
 B_{8,1}=f_2^8\frac{f_8\circ t_{-G}}{f_8},
\]
where $f_2,f_4,f_8$ are the chronological original tangent quotients
of (O45.18).  The single nonzero constant $C_*$ is determined by
evaluating the prescribed initial map at a protected regular point
in the same source and target Cartier frames, as in (O45.20), and
is then kept fixed throughout this source-parameter continuation.

The ordered selected-node source cover and its literal source map are
\[
\overline{\mathcal C}:\ v^2=4u^3+3,\qquad
a=-\frac u3-\frac{u^4}{9},\qquad
t_c=\frac1{2u}-\frac{u^2}{3},\qquad
\xi=\frac1{4u^2}+\frac u3.
\]
On the selected normalization the marked branches, section and inverse
torus coordinate are
\[
x=S^2-2\xi,\qquad y=S(S^2-3\xi),\qquad
c=\frac v{2u},\qquad
Z=\frac{S-c}{S+c},\qquad
P_G=\frac{\sqrt3+v}{\sqrt3-v}.
\]
The open family is obtained by removing exactly the 48 points specified
in O45.2. Its fundamental group has the full marked presentation
\[
[A,B]\mu_1\cdots\mu_{48}=1,
\]
with the two handle loops and the individually specified meridians.

There are exactly twelve points of \(\overline{\mathcal C}\) with
\[
16u^6+120u^3+243=0.
\]
At each of them two unselected nodal fibres combine into a type-II
fibre, while the selected fibre remains nodal, its branch order remains
marked, and the original cap value remains ordinary. On a sufficiently
small disc about each point, the complete original source, \(M_8\)
linearization, prescribed selected fan, infinity attachment and
ordinary cap extend to a proper holomorphic submersion for a fixed
common strict-contraction scalar. Its boundary objects are the original
OS45 objects through the full coefficient maps and their inverses.
The original six-sphere topology is carried by the explicit path family
and the resulting equivariant horizontal flow.

Each of the twelve marked unselected-II meridians has physical return
isotopic to the identity, and hence physical return class \(0\).
For every scalar admitted on the original boundary family, the
constant-argument scalar interval and the full cap logarithm give the
actual annular comparison with this disc family, so the same conclusion
holds at that scalar. Every source, point-line, finite character,
period integer and endpoint marking carried by the original
multiplicity-one construction is retained. The other marked generators
retain their explicit source labels in O45.2; their return classes are
outside this result.
\end{theorem}
\begin{proof}

% Expanded verified cp36 input: notes/part3_os45_unselected_ii_proof.tex
\par\bigskip\noindent\textbf{Complete supporting proof 1: O45.}\par


% Expanded verified cp36 input: notes/part3_os45_unselected_ii_o45_proof.tex
The three supporting arguments retain the original source,
Cartier modules, infinity period words and ordinary-cap reference.
They establish the joint disc family and its compactness,
then transfer its identity return through the actual scalar
and selected-fan morphisms.

\par

\medskip\noindent\textbf{The OS45 marked-source cover and the unselected-II disc attachments}

\par\nobreak

This proof retains the literal source family used in K10 and its
original multiplicity-one cap at the ordinary value \(t=0\). The line
is \(M_8=\mathcal O([8G]-[O])\), the suspension source map is
translation by \(-G\), and the simple zero is the whole selected
\(I_1\) fibre. The scalar-loop calculation is already covered by
K73. Here the parameter moves the marked source. All source,
point-line, selected-height and cap arrows below are specified on
that family. Historical K10 return values are not used.

\par

\medskip\noindent\textbf{O45.1. The literal source and its marked-node cover}

\par\nobreak

Write
\[
\begin{gathered}
 R=t^2+a,\qquad f=R^2+bt,\qquad
 g=R^3+\frac32btR+\frac38b^2,\\
 S_{a,b}:\ y^2=x^3-3fx+2g,\qquad
 G=(R,\sqrt3 b/2).
 \end{gathered}
\tag{O45.1}
\]
The fixed square root is the one in K10. The following scaling is an
actual map of the source, section and relative invariant differential:
\[
 (t,x,y)=(sT,s^2X,s^3Y),\qquad
 a=s^2A,\quad b=s^3B,\qquad
 \frac{dx}{y}=s^{-1}\frac{dX}{Y}.
 \tag{O45.2}
\]
The reciprocal substitutions are its inverse. We work on the
explicit \(B=1\) slice after adjoining the specified cube root of
\(b\ne0\). Thus no scalar or source-coordinate identification is
being made without this map.

Set \(D=64(g^2-f^3)/b^2\). Direct expansion and elimination give
\[
\begin{split}
 D(t;a,b)&=48a(t^2+a)^2+8bt^3+72abt+9b^2,\\
 \operatorname{disc}_tD&=331776b^2(256a^3-3b^2)^3,\\
 D(0;a,b)&=3(16a^3+3b^2).
 \end{split}
\tag{O45.3}
\]
The Weierstrass discriminant is \(-27b^2D\). For \(a\ne0\),
the minimal infinity coefficients are units and its order is eight.
Away from \(256a^3=3b^2\), the other four fibres are simple \(I_1\)\textquotesingle s.
At a finite repeated discriminant value, the equations below give
the actual cusp rather than a double nodal total-space singularity.

For a selected double cubic root \(x=\xi\), put
\(u=\xi-(t_c^2+a)\). The equations \(f=\xi^2\), \(g=\xi^3\)
are equivalent, with \(b,u\ne0\), to
\[
 a=-\frac u3-\frac{u^4}{9b^2},\qquad
 t_c=\frac b{2u}-\frac{u^2}{3b},\qquad
 \xi=\frac{b^2}{4u^2}+\frac u3.
 \tag{O45.4}
\]
To see the inverse directly, translate \(x=R+u\) in the cubic.
The double-root equations become
\(bt=u(u+2R)\) and \(u^2(2u+3R)=3b^2/4\). They yield
(O45.4), while \(u=\xi-(t_c^2+a)\) recovers the original marked
point. This is the actual four-sheeted selected-node cover of
the \(a\)-line on \(b=1\), with no discriminant root renamed.
Its derivatives and selected simplicity are
\[
 \frac{da}{du}=-\frac{4u^3+3b^2}{9b^2},\qquad
 \frac{dt_c}{du}=-\frac{4u^3+3b^2}{6bu^2},\qquad
 D_t(t_c)=-\frac{2(4u^3+3b^2)^3}{27b^3u^2}.
 \tag{O45.5}
\]

The entire selected nodal normalization is
\[
 x=S^2-2\xi,\qquad y=S(S^2-3\xi),\qquad
 S_G=-\frac{\sqrt3 b}{2u}.
 \tag{O45.6}
\]
Adjoin \(v\) with \(v^2=4u^3+3b^2\), and set \(c=v/(2u)\).
Then \(c^2=3\xi\), the two marked branches are \(S=\pm c\),
and the actual zero-section-normalized torus coordinate is
\[
 Z=\frac{S-c}{S+c},\qquad
 P_G=\frac{\sqrt3 b+v}{\sqrt3 b-v}.
 \tag{O45.7}
\]
It sends \(O\) to one and \(G\) to the displayed point. Its
inverse is \(S=c(1+Z)/(1-Z)\). Branch reversal \(v\mapsto-v\)
therefore gives \(Z\mapsto Z^{-1}\), by these literal maps.
No same-fan endpoint is being assigned to that reversal here.

On \(b=1\) the compact ordered cover is the smooth elliptic curve
\(\overline{\mathcal C}:v^2=4u^3+3\). Its residual cube deck is
\[
 (u,v)\longmapsto(\zeta_3u,v),\quad
 (a,t,x,y)\longmapsto(\zeta_3a,\zeta_3^2t,\zeta_3x,y).
 \tag{O45.8}
\]
Substitution in (O45.1) proves this source isomorphism and its
reciprocal. The quotient coordinate is \(v\); its source equation
is still used through the full cubic cover and this endpoint map.

\par

\medskip\noindent\textbf{O45.2. Every boundary point on this ordered cover}

\par\nobreak

Pulling back (O45.3) gives the exact factorization
\[
 256a^3-3b^2
 =-\frac{(4u^3+3b^2)^2
 (16u^6+120b^2u^3+243b^4)}{729b^6}.
 \tag{O45.9}
\]
The discriminant of the three unselected roots, after dividing
the actual polynomial \(D\) by \(t-t_c\), is
\[
 -\frac{1024u^4(16u^6+120b^2u^3+243b^4)^3}{6561b^{10}}.
 \tag{O45.10}
\]
Thus on \(b=1\) the following list gives the complete removed
points for the open source family with four simple nodes,
ordered selected branches and ordinary cap \(t=0\).

\begin{longtable}[]{@{}
  >{\raggedright\arraybackslash}p{(\columnwidth - 4\tabcolsep) * \real{0.3000}}
  >{\raggedleft\arraybackslash}p{(\columnwidth - 4\tabcolsep) * \real{0.4000}}
  >{\raggedright\arraybackslash}p{(\columnwidth - 4\tabcolsep) * \real{0.3000}}@{}}
\toprule\noalign{}
\begin{minipage}[b]{\linewidth}\raggedright
Exact locus on \(\overline{\mathcal C}\)
\end{minipage} & \begin{minipage}[b]{\linewidth}\raggedleft
Number of points
\end{minipage} & \begin{minipage}[b]{\linewidth}\raggedright
Source/marking event
\end{minipage} \\
\midrule\noalign{}
\endhead
\bottomrule\noalign{}
\endlastfoot
\(u=0, v=\pm\sqrt3\) & 2 & The selected node approaches infinity at \(a=0\). \\
\(u^3=-3, v=\pm3i\) & 6 & The selected finite node is retained while the infinity fibre becomes \(I_9\). \\
\(4u^3+3=0, v=0\) & 3 & The selected cubic becomes cuspidal. \\
\(16u^6+120u^3+243=0\) & 12 & Two unselected nodes combine into a type-II fibre. \\
\(2187-16u^3(u^3+3)^3=0\) & 24 & The original ordinary cap value \(t=0\) becomes singular. \\
The point at infinity & 1 & The normalized source parameter leaves the finite \(b=1\) chart. \\
\end{longtable}

All rows are disjoint, and every polynomial in the finite rows
has simple roots after the indicated branch multiplicities.
The cap row factors as
\[
 2187-16w(w+3)^3
 =-(2w-3)(8w^3+84w^2+342w+729),\qquad w=u^3.
 \tag{O45.11}
\]
Its \(2u^3=3\) part is exactly \(t_c=0\). The other part retains
the selected node and places an unselected node at the cap.

In the \(v\)-quotient the removed set consists of
\[
 \begin{gathered}
 v=\pm\sqrt3,\quad\infty,\quad\pm3i,\quad0,\\
 v^4+24v^2+162=0,\\
 (v^2-9)(v^6+33v^4+459v^2+4131)=0.
 \end{gathered}
\tag{O45.12}
\]
There are 18 points there. The three points of ramification
of (O45.8) are already removed; the degree-three map between
these open curves is unramified. Its Euler identity is
\(-48=3(2-18)\). On the ordered cover a complete source-loop
system consists of two handle loops and meridians at the 48
listed points, with the single relation
\[
 [A,B]\,\mu_1\cdots\mu_{48}=1.
 \tag{O45.13}
\]
This is an explicit surface presentation, hence 49 free
generators after eliminating one meridian. A cut polygon for
the double cover of the \(u\)-sphere, branched at the three
roots of \(4u^3+3\) and infinity, supplies \(A,B\); paths to
the specified finite algebraic points supply the meridians.
Their markings are transported by (O45.4),(O45.6),(O45.7).
The twelve unselected-II meridians are filled below. The
selected-cusp, \(I_9\), ordinary-cap and remaining cover
generators stay individually labelled; this proof assigns
no value to them.

\par

\medskip\noindent\textbf{O45.3. The actual unselected-II source disc and translation}

\par\nobreak

Fix one point \((u_*,v_*)\) in the unselected-II row. Choose
a small closed analytic disc \(\Delta\) around it avoiding
all other rows. Since \(da/du\ne0\) and \(v_*\ne0\), the
full source and ordered selected branch vary holomorphically
on this disc. Write \(a_*=a(u_*)\), \(c_*=-4a_*\). Then
\[
 c_*^3=-3/4,\qquad a_*=-c_*/4,\qquad
 t_{\mathrm{II}}=-c_*^2=-16a_*^2,\qquad R_{\mathrm{II}}=-c_*.
 \tag{O45.14}
\]
At this point \(f=g=0\). The full coefficient Jacobian is
\[
 \det\frac{\partial(f,g)}{\partial(a,t)}
 =\frac32c_*^2\ne0,\qquad
 \partial_t(x^3-3fx+2g)|_{x=0}=\frac32c_*\ne0.
 \tag{O45.15}
\]
Consequently \((a,x,y)\) are actual local coordinates on the
source total space, with \(t=t(a,x,y)\) the unique holomorphic
solution of its full displayed equation. In particular the
source surface and its projection to \(\Delta\) are smooth
through the cusp. The value \(D(0)=153/16\) retains the
ordinary cap. The selected value is separate because
\[
 t_c+16a^2=\frac{(2u^3+9)(4u^3+3)^2}{162u}
 \tag{O45.16}
\]
is nonzero at every unselected-II point.

The actual translation \(T=-G\), on its affine chord chart,
is
\[
 L=\frac{y+\sqrt3/2}{x-R},\qquad
 x_T=L^2-x-R,\qquad
 y_T=L(x-x_T)-y.
 \tag{O45.17}
\]
At the target the plus-\(G\) slope is \(-L\); its full chord
formula returns \(x,y\). At the cusp \(x-R=c_*\ne0\), so
both (O45.17) and its inverse are holomorphic. They fix the
cusp. In the \((x,y)\) source coordinates their derivative is
\(\left(\begin{smallmatrix}1&2g_G\\0&1\end{smallmatrix}\right)\),
where \(g_G=x(G)/y(G)=-2c_*/\sqrt3\). Thus this is a smooth
source morphism through the complete cusp, with its full
source-parameter derivatives retained. Away from this chart
the usual point/O addition charts and their actual reciprocal
chords apply. At infinity the fixed \(I_8\) simultaneous
resolution and its section marking are retained, as specified
in O45.5 below.

\par

\medskip\noindent\textbf{O45.4. The complete \(M_8\) coefficient, including the cusp module}

\par\nobreak

Use the actual doubling points \(2G,4G,8G\) from (O45.1).
For \(j=1,2,4\), write \(\ell_{jG}\) for the tangent line
at \(jG\), evaluated at the source point, and put
\[
 f_2=\frac{\ell_G}{x-x(2G)},\qquad
 f_4=f_2^2\frac{\ell_{2G}}{x-x(4G)},\qquad
 f_8=f_4^2\frac{\ell_{4G}}{x-x(8G)}.
 \tag{O45.18}
\]
These are complete rational functions, with every addition
coefficient taken from (O45.1). Their horizontal divisors
are \(m[G]-[mG]-(m-1)[O]\). Define
\[
 B_{8,1}(p)=f_2(p)^8\frac{f_8(Tp)}{f_8(p)}.
 \tag{O45.19}
\]
For a meromorphic point section \(\sigma_8\), with divisor
\([8G]-[O]\), its directed line map is
\(\sigma_8(p)\mapsto B_{8,1}(p)\sigma_8(Tp)\).
Indeed its horizontal divisor is exactly
\([G]+[8G]-[9G]-[O]=D-T^*D\), \(D=[8G]-[O]\).
Thus it is the actual Hom map \(M_8\to T^*M_8\).
It is also the pullback by \(T\) of the original
\(\operatorname{Hom}(t_G^*M_8,M_8)\) linearization: its
evaluation point is \(p-G\). This specifies the direction.

No finite whole-fibre coefficient is suppressed in (O45.19).
At each finite irreducible fibre the orders of \(f_8\circ T\)
and \(f_8\) are equal. The denominator of the \(G\)-doubling
has constant \(2y(G)=\sqrt3\), so \(f_2\) has no finite
whole-fibre zero or pole. The actual point and O generators
cancel every remaining horizontal coefficient, including
section coincidences. At \(I_8\), \(8G\) and \(O\) meet the
identity component, so the Hom has multidegree zero. The
cyclic intersection matrix has kernel generated by the full
fibre. Its degree is one: in the present direction the height
formula gives \(G\cdot O=0\), \(8G\cdot O=3\),
\(9G\cdot O=5\), hence
\(\deg\operatorname{Hom}(M_8,T^*M_8)=5-(3+1)=1\).
Equivalently the original direction uses
\(8G\cdot O+1-7G\cdot O+G\cdot O=3+1-3=1\).
The kernel and degree identify its full vertical divisor as
the single infinity fibre. Therefore the literal map
\[
 B_\lambda(p,v_M)=
 \bigl(Tp,\lambda C_* (t-t_c)B_{8,1}(p)v_M\bigr)
 \tag{O45.20}
\]
extends across infinity and has precisely the selected
whole-fibre zero. Here is the definition of the continuation
family used in this proof. Let \(a_0\) denote its initial marked
source and \(L_{a_0}\) the prescribed directed whole-zero line map.
At a protected regular point \(s_0\), evaluate both maps in the
same Cartier frames and set
\[
 C_*=
 \frac{L_{a_0}(s_0)}
 {(t(s_0)-t_c(a_0))B_{8,1,a_0}(s_0)}\ne0.
\]
The ratio is an invertible function on the compact initial source,
so it is constant on that source. Thus (O45.20) agrees with
\(\lambda L_{a_0}\) on the entire initial source. For every
subsequent marked parameter, the continuation family is defined
by (O45.20) with this fixed \(C_*\). The full coefficient maps
and inverses establish the analytic extension of this defined
family. Its boundary family is its restriction with the prescribed
initial cap and reference continued by the displayed maps.

The initial evaluation proves equality at \(a_0\); it does not
identify an independently prescribed varying linearization
\(L_a\) with the defined continuation. Its exact fibrewise ratio,
evaluated in the same frames on a protected section \(s(a)\), is
\[
 c(a)=\frac{L_a(s(a))}
 {(t(s(a))-t_c(a))B_{8,1,a}(s(a))}.
\]
Such an additional comparison retains \(c(a)\) in the full
suspension coefficient. No constancy or disc extension of that
additional ratio is deduced from the one initial evaluation.

Here is the local cusp coefficient in full, beyond its divisor.
On the central cuspidal curve put \(p=x/y\), so
\(x=p^{-2},y=p^{-3}\), and \(G\) has additive parameter
\(g_G\ne0\). Its multiples have parameters \(mg_G\), and
the literal source map is \(p\mapsto p-g_G\). The full
addition formulas (O45.18) give
\[
 f_m(p)=\frac{(-1)^{m-1}m}{g_G^{m-1}}
 \frac{(p-g_G)^m}{(p-mg_G)p^{m-1}},\quad m=2,4,8,
 \qquad
 B_{8,1}(p)=\frac{256}{g_G^8}
 \frac{(p-g_G)(p-8g_G)}{p(p-9g_G)}.
 \tag{O45.21}
\]
On the finite normalization chart the point/O generator ratio
is \(r(p)=(p-8g_G)/p\). Hence its genuine coefficient is
\[
 B_{8,1}(p)\frac{r(p-g_G)}{r(p)}=\frac{256}{g_G^8}.
 \tag{O45.22}
\]
This chart includes the two strict section directions through
their actual Cartier generators. At the cusp itself the frame
is \(\sigma_8\), because neither horizontal divisor passes
through it. Its coefficient, in the actual cusp local ring,
is
\[
 B_{8,1}(x,y)=\frac{256}{g_G^8}
 \frac{1-73g_G^2x+72g_G^3y}{1-81g_G^2x}.
 \tag{O45.23}
\]
Both numerator and denominator are units at \((0,0)\).
The overlap between these two frames is the full \(r(p)\).
In particular the point/O frame is not being extended across
the cusp by deleting its linear normalization coefficient.
Equations (O45.22),(O45.23) give the precise gluing instead.
The inverse coefficient is evaluated at the actual inverse
source point. Thus the source, point line and entire Hom map
extend on this disc with no charged or deleted cusp point.

\par

\medskip\noindent\textbf{O45.5. Both retained collars and the original ordinary cap}

\par\nobreak

At the selected value put \(q=t-t_c(u)\), preserving this
original height. Let \(\xi(q,u)=\sqrt{f(t_c+q,u)}\) be the
branch with (O45.4) at \(q=0\). It is a unit throughout the
small selected collar. Write

\[
 E(q,u)=g(t_c+q,u)-\xi(q,u)^3,\quad
 h(q,u)=2E(q,u)/q.
 \tag{O45.24}
\]
Here \(h(0,u)=-2u^4\), hence \(h\) is an actual extending
unit. With the ordered square root of \(x+2\xi\), put
\[
 X=y-(x-\xi)\sqrt{x+2\xi},\qquad
 Y=\frac{y+(x-\xi)\sqrt{x+2\xi}}{h(q,u)}.
 \tag{O45.25}
\]
Then \(XY=q\) literally. Conversely, \(q=XY\),
\(y=(X+hY)/2\), and the implicit equation
\(2(x-\xi)\sqrt{x+2\xi}=hY-X\) has derivative
\(2\sqrt{3\xi}\ne0\), giving its unique holomorphic
inverse. Thus the actual nodal source chart and every unit
are specified before applying the cusp fan.

Let \(Q=\kappa(q,u)q\) be the actual Tate parameter from
the original \(j\)-function, with \(Qj(Q)=1+O(Q)\), and
let \(Z\) be the inverse marked uniformization agreeing
with (O45.7). Define \(P=Z(G)\) on the full smooth collar.
Both its differential scale and the inverse uniformization
are chosen to agree with the original reference; their
unique analytic continuations over \(\Delta\) retain all
higher coefficients. For the complete theta product
\(f_b(Z)=\Theta_Q(Z/b)/\Theta_Q(Z)\), the point coefficient
\(V_{\rm pt}=f_{P^8}(Z)v_M\) has exact periods
\[
 h:(Z,V_{\rm pt})\mapsto(QZ,P^8V_{\rm pt}),\qquad
 B:(Z,V_{\rm pt})\mapsto(P^{-1}Z,\lambda q\beta V_{\rm pt}),
 \tag{O45.26}
\]
where, as a point-module map,
\[
 \beta=C_* B_{8,1}(p)
 \frac{f_{P^8}(P^{-1}Z)}{f_{P^8}(Z)}.
 \tag{O45.27}
\]
The horizontal divisors cancel in this displayed direction,
so (O45.27) is independent of the point \(p\). Evaluation
at one protected regular \(Z_*\) gives its full analytic
coefficient and inverse. It is a unit by (O45.20). The
usual selected toric charts, their prescribed fan and every
period translate therefore vary holomorphically over the
entire disc. Their source height is the same \(q\); the
units \(\kappa,P,\beta\) have not been discarded.

At infinity retain \(z=1/t\) and the minimal coordinates
\((z^2x,z^3y)\). The exact source has a split node with
smoothing parameter \(z^8\kappa_\infty(z,u)\), where
\(\kappa_\infty(0,u)=-a(u)/2304\). Choose its branch
so \(P_\infty=z^3\pi(z,u)\), \(\pi(0,u)=-1/24\).
These functions are the full inverse Tate uniformization,
not their leading terms. The source coordinates have their
actual consecutive charts \(X_iY_i=z\), obtained from
\(X_i=Z/z^i,Y_i=z^{i+1}/Z\), with the corresponding full
source units. Their overlaps and translation are exactly
the monomial arrows obtained by the displayed \(Q,P\).
The section remains on component three; its inverse
translation remains the shift by minus three.

The actual identity-component point of \(8G\) is
\[
 b_0=P_\infty^8/Q_\infty^3,\qquad
 b_0(0,u)=-1/(9a(u)^3).
 \tag{O45.28}
\]
Three original group-law doublings verify this value, as
well as the full minimal coordinates
\[
 X_{8,0}=\frac{81a^6-90a^3+1}{(9a^3+1)^2},\quad
 (Y_{8,0}/\sqrt3)=
 \frac{108a^3(9a^3-1)}{(9a^3+1)^3}.
 \tag{O45.29}
\]
At the II centre \(b_0=-256/27\ne1\), so both point/O
modules are separated on this infinity collar. The full
theta identity is
\[
 f_{Q^3b_0}(Z)=(-Z/b_0)^3Q^{-6}f_{b_0}(Z).
 \tag{O45.30}
\]
It gives the actual reciprocal line-frame arrow, including
the base power \(Q^{-6}\). In the \(b_0\) point frame the
suspension coefficient is \(P_\infty^3\) times its original
point-theta coefficient, and the source/line valuation
vectors are exactly \((8,0),(-3,0)\). The primitive periods
are the full words
\[
 H=h^{-1}B^{-3},\qquad K=h^3B^8,\qquad
 h=H^8K^3,\quad B=H^{-3}K^{-1}.
 \tag{O45.31}
\]
Their integer matrix has determinant one. The line
coefficient of \(K\) is \(\lambda^8\) times its complete
analytic unit. Thus the original nonzero-line \(I_8\)
quotient and every source component persist on the entire
disc, with literal source and line inverses. This is the
same rank-one filling as the source line complement:
on each consecutive source chart its map is the above
identity on the actual point module and the indicated
integer period reindex. Their equations and overlaps
agree, so the maps glue across the whole cycle.

Finally choose the fixed original cap disc around \(t=0\),
disjoint from the selected and unselected source collars.
All its source fibres are smooth over \(\Delta\). Continue
the prescribed full periods and their initial logarithms
over this contractible parameter disc. In particular
\(\gamma(t,u,\lambda)\) is the actual suspension period of
(O45.20), with its full source, line and scalar components.
The original cap-to-ordinary map and inverse are
\[
 \Psi(\zeta,t,u)=
 \left(\zeta+\frac{\log t}{2\pi i}\gamma(t,u,\lambda),t,u\right),
 \qquad
 \Psi^{-1}(\zeta,t,u)=
 \left(\zeta-\frac{\log t}{2\pi i}\gamma(t,u,\lambda),t,u\right).
 \tag{O45.32}
\]
A logarithm branch change adds exactly the original
\(\gamma\)-period. Thus this is the unchanged order-one
cap with its marked meridian and physical line, now
holomorphic over \(\Delta\). All initial annular
integers are carried by analytic continuation from the
reference. There is no source-parameter deck around the
filled disc: every displayed source, module, theta and
cap map returns to itself.

\par

\medskip\noindent\textbf{O45.6. Compact disc family and the filled source meridians}

\par\nobreak

The punctured ordered curve of O45.2 is connected. Choose
a path in it from one of the original K10 \(a=b=1\)
markings to the chosen boundary point of \(\Delta\).
Cover that path by finitely many small discs avoiding
all rows of the table. On each disc the same original
equation, source inverses, simultaneous \(I_8\) charts,
selected fan and cap map (O45.32) apply. On overlaps
the rational point line and (O45.20) are literally
the same; source period changes are their actual
integral changes of basis, with the initial cap
period and reference continued along the path.
A smaller common scalar and its constant-argument
interval identify the original admitted boundary
family with these proper submersions. This carries
the known K10 sphere to the boundary of the present
disc by explicit source/line/cap families; no K10
return value enters this transport.

Choose a Hermitian metric on the actual point line over
the compact source family outside the selected tube.
Equations (O45.17)-\/-(O45.23) give a uniform finite norm
for the full unscaled lift, including the II centre.
Choose the fixed scalar with norm times this bound less
than one. The full suspension is then strictly
contracting there. The quotient is free: a physical
circle equality with \(B^k\) forces \(k=0\) by the norm
inequality. A compact annular fundamental domain gives
properness and Hausdorffness uniformly over the closed
source parameter disc.

Use the actual selected fan charts (O45.25)-\/-(O45.27)
inside the excluded tube, the unchanged infinity
charts (O45.28)-\/-(O45.31), and the original cap
(O45.32). Their overlaps are the full source, point
module and period arrows already written. There are
finitely many chart and period representatives on
each compact collar. They glue to a proper holomorphic
family over \(\Delta\). The source projection is a
submersion at the unselected cusp by (O45.15), at the
selected fan by its actual toric coordinates, and at
the other pieces by their unchanged source or cap
coordinates. Hence the completed family is a smooth
proper submersion over the whole disc.

On its punctured boundary these are exactly the
original OS45 objects with the selected whole \(I_1\),
\(M_8\), translation \(-G\), fixed fan and original
ordinary cap, by the explicit line-reference match
in (O45.20). The original sphere fibres transport
through this actual proper submersion. A connection
on this disc family, averaged over the physical
circle, retains its line action and provides the
disc trivialization. Its boundary return is isotopic
to the identity. Thus each of the twelve marked
unselected-II meridians has zero physical return;
this follows from its displayed filling, with all
source and cap markings retained.

The same conclusion holds for every scalar admitted
on the original boundary family. To compare it with
the common disc-contraction scalar, decrease its
modulus along the actual constant-argument interval.
On the boundary all intermediate lifts remain
contracting. Carry the full \(\log\lambda\) in
\(\gamma\), (O45.27), (O45.31), and (O45.32) along
that interval. This gives the exact relative
parameter family and inverse, with no extra scalar
turn. The smaller endpoint fills by the disc just
constructed, and the boundary annulus transports
that filling to the original admitted scalar.

This closes the unselected-II source generators in
the ordered cover. The selected-cusp, \(I_9\),
ordinary-cap-collision and remaining source-cover
generators of (O45.13) remain outside this completed
disc calculation. Their precise source equations
and parameter labels have been retained above.

\par

\medskip\noindent\textbf{O45.7. Exact source and integral computations}

\par\nobreak

Exact rational arithmetic in
\path{code/current_audits/audit_paper_i/iv_os45/code/k116_afterreading_exact.py}
starts from the displayed source equation and the five literal
K10 matrices. Its 249 checks include the rational source and
translation inverses, chronological cusp Miller functions,
both point/O frames, original doubling at infinity,
the Hom-degree and Gram calculations, primitive period words,
a twelve-coset modular enumeration, and the integral
parabolic complexes for all four selected-zero choices.
The full matrices and results are in
\path{code/current_audits/audit_paper_i/iv_os45/output/k116_afterreading_exact.json}.
The joint chart maps and compactness argument above supply
the analytic disc extension and its return.

% End expanded verified cp36 input


\par\bigskip\noindent\textbf{Complete supporting proof 2: O45R.}\par


% Expanded verified cp36 input: notes/part3_os45_unselected_ii_o45r_proof.tex
\par

\medskip\noindent\textbf{Source-disc proof of the OS45 unselected-II attachments}

\par\nobreak

We retain the whole original source, point line, selected height,
ordinary cap and fixed scalar of O45.1-\/-32. In particular
\[
 R=t^2+a,\quad f=R^2+t,\quad
 g=R^3+\frac32tR+\frac38,\quad
 y^2=x^3-3fx+2g,\quad G=(R,\sqrt3/2).
 \tag{O45R.1}
\]
The chosen square root is fixed. The point line is
\(M_8=\mathcal O([8G]-[O])\), the source translation is
\(T:p\mapsto p-G\), and the ordinary logarithmic cap
at \(t=0\) has its full inward suspension clutch.
The initial source is the literal K10 construction.
Its old scalar-return sign is not used in this source-disc proof.

\par

\medskip\noindent\textbf{O45R.1. Original parameter maps and the exact collision discs}

\par\nobreak

The weighted map
\((t,x,y)=(sT,s^2X,s^3Y)\), \(a=s^2A,b=s^3B\)
has inverse given by division by those powers. Substitution
in the equation multiplies it by \(s^6\); it carries the
actual section and scales \(dx/y\) by \(s^{-1}\).
Thus the \(b=1\) slice is taken on the explicitly marked
cube-root cover of \(b\ne0\). We keep that cover.

Direct differentiation of the original cubic at a selected
double root \(\xi\), followed by putting
\(u=\xi-(t_c^2+a)\), gives
\[
 t_c=u(u+2R_c),\qquad
 u^2(2u+3R_c)=\frac34,\qquad
 R_c=t_c^2+a .
 \tag{O45R.2}
\]
Solving these two equations gives
\[
 a=-u/3-u^4/9,\quad
 t_c=1/(2u)-u^2/3,\quad
 \xi=1/(4u^2)+u/3.
 \tag{O45R.3}
\]
Their inverse is the original difference \(u=\xi-R_c\).
The polynomial
\[
 D=48a(t^2+a)^2+8t^3+72at+9
 \tag{O45R.4}
\]
satisfies \(\Delta=-27D\). Eliminating \(t\) gives
\(\operatorname{disc}_tD=331776(256a^3-3)^3\).
The selected-root derivative is
\[
 D_t(t_c)=-\frac{2(4u^3+3)^3}{27u^2}.
 \tag{O45R.5}
\]
Consequently the branch cover is the literal elliptic curve
\(v^2=4u^3+3\), with the excluded loci specified in O45.2.
On its selected fibre the maps
\[
 x=S^2-2\xi,\quad y=S(S^2-3\xi),\qquad
 Z=\frac{S-v/(2u)}{S+v/(2u)}
 \tag{O45R.6}
\]
have inverses \(S=y/(x-\xi)\) and
\(S=(v/(2u))(1+Z)/(1-Z)\) on their indicated charts.
They send the original \(G\) to
\((\sqrt3+v)/(\sqrt3-v)\). Inversion of the branch marking
is literally \(Z\mapsto Z^{-1}\). We do not use it as a
closing of a fixed fan.

The unselected collision polynomial is
\[
 H(u)=16u^6+120u^3+243.
 \tag{O45R.7}
\]
It has six simple roots, disjoint from \(u=0\),
\(u^3=-3\), \(4u^3+3=0\) and
\(2187-16u^3(u^3+3)^3=0\). Each has two distinct
values of \(v\), so there are precisely twelve
marked disc centres. These are all the unselected
collisions: division of O45R.4 by its actual factor
\(t-t_c\) gives cubic discriminant
\(-1024u^4H(u)^3/6561\). The distinct selected
root and the ordinary cap therefore stay separate.

The other loci in O45.2 have, respectively, two, six,
three, twenty-four and one points. Direct polynomial
gcd calculations show that their rows are disjoint.
Their total with the twelve centres is forty-eight.
The curve \(v^2=4u^3+3\) has four simple branch
points over the \(u\)-sphere, hence genus one by
its actual double-cover Euler calculation. Deleting
the forty-eight points gives a connected surface.
The map to \(v\) is a degree-three cover, with its
three ramification points among the deleted points.
Its eighteen-puncture target has Euler number
\(-16\), and the source has Euler number \(-48\).
The surface presentation with two handle loops and
the forty-eight meridians has the single boundary
relation written in O45.13. The present disc maps
fill exactly the twelve meridians of O45R.7.

Fix one centre and a closed disc \(\Delta\) around it.
The maps in O45R.3 and the chosen branch \(v\) are
holomorphic on \(\Delta\); \(da/du\ne0\) there.
Write \(c=-4a_*\). Then the collision point is
\[
 c^3=-3/4,\quad a_*=-c/4,\quad
 t_*=-c^2,\quad R_*=-c,\quad f_*=g_*=0 .
 \tag{O45R.8}
\]
The independent original-equation derivatives are
\[
 \det\frac{\partial(f,g)}{\partial(a,t)}
    =\frac32c^2,\qquad
 \partial_t(x^3-3fx+2g)|_{x=y=0}
    =\frac32c .
 \tag{O45R.9}
\]
Both are nonzero. The full source near its cusp is
therefore the graph of a holomorphic function
\(t=t(a,x,y)\); projection to \(a\) is a submersion.
This is a chart on the original total source, with
inverse the original \((a,x,y)\) projection.
The cusp occurs at a finite value different from
the selected value, since
\[
 t_c+16a^2=
 \frac{(2u^3+9)(4u^3+3)^2}{162u}\ne0
 \quad\text{on }H(u)=0.
 \tag{O45R.10}
\]
The original cap remains ordinary because its
discriminant factor is \(D(0)=153/16\).

\par

\medskip\noindent\textbf{O45R.2. Full source translation and the two Cartier frames}

\par\nobreak

The actual translation by \(-G\) is
\[
 L=\frac{y+\sqrt3/2}{x-R},\qquad
 x_T=L^2-x-R,\qquad y_T=L(x-x_T)-y .
 \tag{O45R.11}
\]
Substitution reduces its source equation to the
original source equation. Its inverse uses the
plus-\(G\) slope
\((y_T-\sqrt3/2)/(x_T-R)=-L\); applying the
same chord formula gives \(x,y\) exactly.
At the cusp both denominators are \(c\ne0\),
both maps fix \(x=y=0\), and their derivatives
in the graph coordinates of O45R.9 are
\[
 DT=\begin{pmatrix}1&2g_G\\0&1\end{pmatrix},
 \qquad
 DT^{-1}=\begin{pmatrix}1&-2g_G\\0&1\end{pmatrix},
 \qquad g_G=-2c/\sqrt3 .
 \tag{O45R.12}
\]
This proves invertibility on the full cusp chart.
The original section and zero-section charts give
the same inverse on the rest of the source.

Let \(f_m\) be the actual Miller function obtained
by consecutive doublings of \(G\), for \(m=2,4,8\).
Its horizontal divisor is \(m[G]-[mG]-(m-1)[O]\).
The full rational coefficient
\[
 B_{8,1}=f_2^8\,\frac{f_8\circ T}{f_8}
 \tag{O45R.13}
\]
has horizontal divisor
\([G]+[8G]-[9G]-[O]\). The point-line frame
at the target has divisor \(T^*([8G]-[O])
=[9G]-[G]\). Thus multiplication by O45R.13
carries the actual two Cartier modules in the
stated direction. All horizontal zeros and
poles cancel in those modules, at their
intersections as well as away from them.

We reconstruct its central value as a function
before restricting it to a constant. Put
\(p=x/y\) on \(y^2=x^3\), so that
\(x=p^{-2},y=p^{-3}\) and \(G\) has additive
coordinate \(g_G\). The tangent at \(jG\)
has slope \(3/(2jg_G)\). Substituting this
tangent and \(x(2jG)=(2jg_G)^{-2}\) in its
chord quotient gives all three identities
\[
 f_m(p)=\frac{(-1)^{m-1}m}{g_G^{m-1}}\,
       \frac{(p-g_G)^m}{(p-mg_G)p^{m-1}},
 \qquad m=2,4,8 .
 \tag{O45R.14}
\]
Using \(Tp=p-g_G\) in O45R.13 now gives
\[
 B_{8,1}(p)=\frac{256}{g_G^8}
       \frac{(p-g_G)(p-8g_G)}{p(p-9g_G)}.
 \tag{O45R.15}
\]
The finite normalization chart uses the true
point/O ratio \(r(p)=(p-8g_G)/p\).
Its coefficient is
\[
 B_{8,1}(p)\frac{r(p-g_G)}{r(p)}
       =\frac{256}{g_G^8}.
 \tag{O45R.16}
\]
The cusp chart uses the original meromorphic
point frame, whose horizontal divisor avoids
the cusp. Multiplying numerator and denominator
of O45R.15 by the exact opposite factor
\(1+9g_G/p\) gives the literal original-ring expression
\[
 B_{8,1}(x,y)=\frac{256}{g_G^8}
       \frac{1-73g_G^2x+72g_G^3y}
            {1-81g_G^2x}.
 \tag{O45R.17}
\]
Both factors have value one at the cusp.
The overlap of these two frames is the full
function \(r\). In particular O45R.17, and
its reciprocal at the inverse source point,
are actual holomorphic module maps through
the cusp. The parameter graph O45R.9 and the
nonzero denominators in these formulas make
the maps holomorphic on a neighborhood in
the total family.

At every finite irreducible fibre the two
terms \(f_8\circ T\) and \(f_8\) have the same
whole-fibre order. The \(G\)-doubling denominator
is the fixed nonzero number \(\sqrt3\), and
the divisor of \(f_2\) has no finite
whole-fibre coefficient. Thus O45R.13 has
no finite vertical divisor after its
horizontal module cancellation. At infinity
\(8G\) and \(O\) lie on the identity
component, so the Hom has zero component
degrees. The \(I_8\) intersection matrix
has kernel \(\mathbb Z(1,\ldots,1)\).
The height \(h(G)=1/8\) and component
index \(3\) give
\[
 G\cdot O=0,\quad 7G\cdot O=3,\quad
 8G\cdot O=3,\quad 9G\cdot O=5 .
 \tag{O45R.18}
\]
Restriction of the directed Hom to \(O\)
has degree \(5-0-3-1=1\). Its entire
vertical divisor is consequently one
infinity fibre. Multiplication by \(t-t_c\)
replaces this by the selected whole zero,
giving exactly the original suspension
\[
 B_\lambda(p,v_M)=
  (Tp,\lambda C_*(t-t_c)B_{8,1}(p)v_M).
 \tag{O45R.19}
\]
The prescribed nonzero \(C_*\) is fixed by
one original regular-point evaluation.
The matching ratio is an invertible
function on the compact source and is
constant there. Continuing that constant
retains the initial scalar normalization.

\par

\medskip\noindent\textbf{O45R.3. The selected and infinity collars with every unit}

\par\nobreak

At the selected fibre use its original
\(q=t-t_c(u)\). Choose
\(\xi(q,u)=\sqrt{f(t_c+q,u)}\) with the
marked value O45R.3, and put
\[
 h(q,u)=\frac{2(g(t_c+q,u)-\xi(q,u)^3)}q.
 \tag{O45R.20}
\]
The numerator vanishes at \(q=0\), and its
quotient has value \(-2u^4\ne0\). Thus
\[
 X=y-(x-\xi)\sqrt{x+2\xi},\qquad
 Y=h^{-1}\bigl(y+(x-\xi)\sqrt{x+2\xi}\bigr)
 \tag{O45R.21}
\]
has exact equation \(XY=q\). The inverse
first puts \(q=XY\), \(y=(X+hY)/2\),
then solves
\(2(x-\xi)\sqrt{x+2\xi}=hY-X\).
Its derivative at the node is
\(2\sqrt{3\xi}\ne0\). Both directions are
therefore holomorphic on the complete
selected collar over \(\Delta\).

The full inverse Tate uniformization gives
\(Q=q\kappa(q,u)\), the marked \(Z\), and
\(P=Z(G)\), with their original reference
scales. The point line is carried by
\[
 V_{\mathrm{pt}}=
 \frac{\Theta_Q(Z/P^8)}{\Theta_Q(Z)}\,v_M.
 \tag{O45R.22}
\]
Its full period arrows are
\[
 h:(Z,V_{\mathrm{pt}})\mapsto(QZ,P^8V_{\mathrm{pt}}),
 \qquad
 B:(Z,V_{\mathrm{pt}})\mapsto
        (P^{-1}Z,\lambda q\beta V_{\mathrm{pt}}),
 \tag{O45R.23}
\]
where
\[
 \beta=C_* B_{8,1}(p)
       \frac{f_{P^8}(P^{-1}Z)}{f_{P^8}(Z)} .
 \tag{O45R.24}
\]
The displayed Cartier divisors cancel, and
the ratio is independent of \(p\) by
compactness of a regular source fibre.
Evaluation at a protected regular point
defines its complete analytic base function.
The same module divisor calculation and
its reciprocal prove that it is an
extending unit. Thus every coefficient of
the prescribed selected fan is analytic on
\(\Delta\). On an affine cone the transition
acts on each character by its actual
monomial times these full units; the
inverse monomial and reciprocal units
give its inverse. The ordered branch is
single-valued on \(\Delta\), so this
extends the original prescribed fan
without a fan permutation at its centre.

There is a separate check of the infinity
normalization. Put \(z=1/t\) and use the
original minimal coordinates
\(X=z^2x,Y=z^3y\). The leading inverse
\(j\)-coefficient gives
\[
 Q_\infty=z^8\kappa_\infty(z,u),\qquad
 \kappa_\infty(0,u)=-a/2304 .
 \tag{O45R.25}
\]
The central cubic is
\(Y^2=(X-1)^2(X+2)\). Its exact normalized
torus coordinate has
\[
 X=1+\frac{12Z}{(1-Z)^2},\qquad
 \frac{Y}{\sqrt3}
      =\frac{12Z(Z+1)}{(Z-1)^3}.
 \tag{O45R.26}
\]
The zero section has \(Z=1\); this fixes
the coefficient twelve and the branch
sign. For the original section \(G\),
\[
 X_G=1+az^2,\quad
 Y_G=(\sqrt3/2)z^3,\quad
 \sqrt{f_{\min}}=1+az^2+z^3/2+O(z^5).
 \tag{O45R.27}
\]
The marked component is three. The full
Tate expression for its branch has leading
term \(X_G-\sqrt{f_{\min}}=12P_\infty\);
the next Laurent terms have orders at least
five, from \(Q_\infty/P_\infty\), and six,
from \(P_\infty^2\). Its analytic scaling
has central value one by O45R.26.
It follows that
\[
 P_\infty=z^3\pi(z,u),\qquad
 \pi(0,u)=-1/24,\qquad
 \left.\frac{P_\infty^8}{Q_\infty^3}\right|_{z=0}
          =-\frac1{9a^3}.
 \tag{O45R.28}
\]
The full functions, rather than their
leading terms, define the subsequent
maps. Formula O45R.26 applied to the
last value also gives independently
\[
 X_{8,0}=\frac{81a^6-90a^3+1}{(9a^3+1)^2},
 \quad
 \frac{Y_{8,0}}{\sqrt3}
       =\frac{108a^3(9a^3-1)}{(9a^3+1)^3}.
 \tag{O45R.29}
\]
This agrees with the full original
addition-law calculation in O45.29.
At each unselected-II centre the point
value is \(-256/27\), and hence is
distinct from the zero section.

Put \(b_0=P_\infty^8/Q_\infty^3\).
The product definition of theta gives
the entire identity
\[
 f_{Q_\infty^3b_0}(Z)=
       (-Z/b_0)^3Q_\infty^{-6}f_{b_0}(Z).
 \tag{O45R.30}
\]
Its reciprocal is the actual change
from the original point frame to the
identity-component point frame. Under
translation by \(P_\infty^{-1}\), its
factor changes the full line coefficient
by \(P_\infty^3\). This retains the
base factor \(Q_\infty^{-6}\) in the
frame arrow. The valuation vectors
in that frame are \((8,0)\) and
\((-3,0)\).

The whole period reindex is
\[
 H=h^{-1}B^{-3},\quad K=h^3B^8,\qquad
 h=H^8K^3,\quad B=H^{-3}K^{-1}.
 \tag{O45R.31}
\]
Its determinant is one and its
source valuations are \(1,0\).
The second multiplier of \(K\)
is \(\lambda^8\) times its full
extending unit. The consecutive
source charts
\(X_i=Z/z^i,Y_i=z^{i+1}/Z\)
have equation \(X_iY_i=z\).
Their full \(Q_\infty/z^8\) and
\(P_\infty/z^3\) units give all
period and translation overlaps,
with the inverse overlaps given
by their reciprocals. Thus the
rank-one completed quotient and
the original source line
complement are identified chart
by chart, including the entire
eight-component cycle.

\par

\medskip\noindent\textbf{O45R.4. The original cap and the proper disc family}

\par\nobreak

The disc at \(t=0\) is smooth over
\(\Delta\), by \(D(0)=153/16\).
Its original period basis and
suspension logarithm continue
holomorphically on this
contractible parameter disc.
Retain their full vector
\(\gamma(t,u,\lambda)\).
The cap-to-ordinary map is
\[
 \Psi(\zeta,t,u)=
  \left(\zeta+\frac{\log t}{2\pi i}
           \gamma(t,u,\lambda),t,u\right).
 \tag{O45R.32}
\]
The inverse subtracts precisely
this full vector. A logarithm
branch change adds one actual
\(\gamma\)-period, so both maps
are single-valued on the
original torus quotient. This
preserves the original
multiplicity-one cap, its
nonzero meridian clutch and
every initial annular integer.
It also preserves the
physical line: the
coefficient vector contains
its original line coordinate
and the complete
\(\log\lambda\).

All the full source and line
maps above are holomorphic
over the closed parameter
disc, and the selected
coefficient has its sole
zero on the removed selected
fibre. On the complementary
compact source choose a
Hermitian metric on \(M_8\).
The entire unscaled lift has
a uniform finite norm.
A smaller common nonzero
scalar makes that norm
strictly less than one.
The inverse lift on its
image expands. A fundamental
annular shell and its two
identified faces therefore
give a compact proper
Hausdorff quotient, uniformly
over \(\Delta\).

The cusp chart O45R.9 is
smooth over the source
parameter, and O45R.17
extends the line map there.
The selected toric charts
O45R.21-\/-24 are smooth
over that parameter.
The infinity charts
O45R.25-\/-31 and the
ordinary cap O45R.32
are smooth over it as well.
Their displayed inverses
and common source/line
overlaps glue these pieces
to a proper holomorphic
submersion over
\(\Delta\).

The topology is transported
from the original K10
source by an actual path
in the connected punctured
ordered cover. On finitely
many discs along that path
the same source equations,
point modules, selected
fan, infinity charts and
ordinary cap apply.
Their original period
changes are integral
basis changes with their
full source/line
coefficients. A smaller
common scalar gives
proper submersions
over the whole compact
path. A horizontal lift
of its tangent field,
patched on the compact
quotient by a partition
of unity, integrates to
an oriented fibre
diffeomorphism. Its
inverse is the reverse
flow. It carries the
actual K10 six-sphere
and its prescribed
fan to the boundary
of the current disc.

The same construction on
the filled disc gives a
smooth trivialization:
choose a connection,
average over the compact
physical circle, and
transport radially
from the centre.
Its inverse uses the
reverse radial paths.
All endpoints on a
boundary circle are
therefore identified
by a product whose
closing is the
identity. The original
analytic family has
this same closing,
because every source,
point, fan and cap
coefficient above is
single-valued on
\(\Delta\).
Its physical return
is isotopic to the
identity, hence has
class zero.

Finally take any scalar
admitted on the original
boundary family. The
constant-argument radial
interval to the smaller
common scalar stays in
the strict contraction
domain on that boundary.
Substitution of this
interval in the complete
periods, O45R.19,
O45R.24, O45R.31 and
O45R.32 gives an
actual proper annular
comparison family.
It has no scalar
winding and retains
the endpoint maps.
Horizontal flow on
that compact annulus,
with its reverse as
inverse, gives an
isotopy of the two
boundary returns.
Thus every admitted
fixed scalar gives
class zero on each
of the twelve marked
unselected-II
meridians. The other
source-cover
generators remain
outside this
calculation.

% End expanded verified cp36 input


\par\bigskip\noindent\textbf{Complete supporting proof 3: O45P.}\par


% Expanded verified cp36 input: notes/part3_os45_unselected_ii_o45p_proof.tex
\par

\medskip\noindent\textbf{Complete source and cap comparison for the twelve OS45 unselected-II meridians}

\par\nobreak

This proof retains the original source, section, point line, suspension,
ordinary logarithmic cap, selected branch order, and prescribed lattice-basis
fan of O45.1-\/-7. It compares that source-disc construction with the
full source and cap maps of O45R.1-\/-4. Its conclusion is the zero return of the
twelve specified source meridians. The other source generators retain their
separate obligations. The earlier numerical scalar-return sign in K10 is
not an input to this proof: the actual source, point line, local-system
maps, integral topology and ordinary cap of that construction are used.

\par

\medskip\noindent\textbf{O45P.1. Original source maps and the exact marked discs}

\par\nobreak

Write
\[
 R=t^2+a,\quad f=R^2+bt,\quad
 g=R^3+\frac32btR+\frac38b^2,\quad
 y^2=x^3-3fx+2g,\qquad G=(R,\sqrt3b/2).
 \tag{O45P.1}
\]
The chosen square root is fixed. The weighted arrow
\((t,x,y)=(sT,s^2X,s^3Y)\), \(a=s^2A,\ b=s^3B\), and its inverse
divide the equation by \(s^6\) and carry \(dx/y\) to \(s^{-1}dX/Y\).
We use its explicitly marked cube-root cover to put \(b=1\).
This keeps the deck arrow
\((u,v)\mapsto(\zeta_3u,v)\) and the original source map
\((a,t,x,y)\mapsto(\zeta_3a,\zeta_3^2t,\zeta_3x,y)\).

For clarity, rationality and the direction of the translation can be
checked on the original source before choosing any singular fibre.
Put \(Y=y/\sqrt3\) and use the two independent coordinates \(A,B\) with
\[
 v_0=3(AB-a),\quad
 t=\frac{B-A}{2}+\frac{b}{2v_0},\quad
 x=t^2+a+v_0,\quad Y=\frac{v_0(A+B)}2.
 \tag{O45P.2}
\]
The inverse is
\(v_0=x-t^2-a,\ A=Y/v_0-t+b/(2v_0),\
B=Y/v_0+t-b/(2v_0)\).
Substitution gives the original equation \(3Y^2=x^3-3fx+2g\).
Thus these are mutual rational source maps on the stated open sets;
the minimal regular compact source is the rational elliptic surface
used in K10.

The full finite discriminant and its discriminant in the base coordinate
are
\[
 D=48a(t^2+a)^2+8bt^3+72abt+9b^2,\quad
 \Delta=-27b^2D,\quad
 \operatorname{disc}_tD=331776b^2(256a^3-3b^2)^3.
 \tag{O45P.3}
\]
At a selected double root \(\xi\) put \(u=\xi-(t_c^2+a)\).
Solving the two original double-root equations gives the actual inverse
marked cover
\[
 a=-\frac u3-\frac{u^4}{9b^2},\qquad
 t_c=\frac b{2u}-\frac{u^2}{3b},\qquad
 \xi=\frac{b^2}{4u^2}+\frac u3.
 \tag{O45P.4}
\]
The inverse is precisely \(u=\xi-(t_c^2+a)\). In particular
\(f(t_c)=\xi^2,\ g(t_c)=\xi^3\), and
\[
 \frac{da}{du}=-\frac{4u^3+3b^2}{9b^2},\qquad
 D_t(t_c)=-\frac{2(4u^3+3b^2)^3}{27b^3u^2}.
 \tag{O45P.5}
\]
The original selected normalization and its inverse are
\[
 x=S^2-2\xi,\quad y=S(S^2-3\xi),\quad S=\frac y{x-\xi},
 \qquad
 Z=\frac{S-v/(2u)}{S+v/(2u)},\quad
 S=\frac v{2u}\frac{1+Z}{1-Z}.
 \tag{O45P.6}
\]
Here \(v^2=4u^3+3b^2\). Substitution of
\(S_G=-\sqrt3b/(2u)\) sends the original section to
\[
 P=\frac{\sqrt3b+v}{\sqrt3b-v}.
 \tag{O45P.7}
\]
Changing the ordered branch is the displayed source arrow
\(Z\mapsto Z^{-1}\). This proof uses the fixed ordered branch on each
disc and its continued fan; it does not use that arrow as a new closing.

Now put \(b=1\) and \(H(u)=16u^6+120u^3+243\). The two exact
factorizations are
\[
 256a(u)^3-3=-\frac{(4u^3+3)^2H(u)}{729},\qquad
 \operatorname{disc}_t\!\left(\frac{D}{t-t_c}\right)
        =-\frac{1024u^4H(u)^3}{6561}.
 \tag{O45P.8}
\]
The six roots of \(H\) are simple. Polynomial gcds with
\(u,\ u^3+3,\ 4u^3+3,\ 2187-16u^3(u^3+3)^3\) are one.
Thus each root has two distinct ordered values of \(v\), and gives one
of the twelve specified unselected-II disc centres.
All forty-eight deleted points are retained: respectively two over
\(u=0\), six over \(u^3=-3\), three with \(v=0\), twelve over \(H=0\),
twenty-four over the ordinary-cap polynomial, and the point at infinity.
The curve \(v^2=4u^3+3\) has four simple branch points over the
\(u\)-sphere; the Euler calculation gives genus one. Removing the
forty-eight points leaves a connected surface with free fundamental
group of rank \(2+48-1=49\). The degree-three map to \(v\) is unramified
on this open surface, and its Euler identity is
\(-48=3(2-18)\). These actual covering maps specify which twelve
meridians are filled.

\par

\medskip\noindent\textbf{O45P.2. Smooth original cusp source and both translation directions}

\par\nobreak

Fix one of the twelve points, and choose a closed parameter disc
\(\Delta\) avoiding all other deleted values. The functions in
O45P.4 and the chosen \(v\) are single-valued on this disc.
At its centre write \(c=-4a_*\). Then
\[
 c^3=-\frac34,\quad a_*=-c/4,\quad
 t_*=-c^2,\quad R_*=-c,\quad f_*=g_*=0.
 \tag{O45P.9}
\]
Direct differentiation of the full coefficients gives
\[
 \det\frac{\partial(f,g)}{\partial(a,t)}=\frac32c^2,\qquad
 \left.\partial_t(x^3-3fx+2g)\right|_{x=y=0}=\frac32c.
 \tag{O45P.10}
\]
Consequently the original source is locally the holomorphic graph
\(t=t(a,x,y)\). Its inverse chart is the original \((a,x,y)\)
projection, so its projection to \(\Delta\) is a submersion.
The centre is disjoint from the selected fibre and the ordinary cap:
\[
 t_c+16a^2=\frac{(2u^3+9)(4u^3+3)^2}{162u}\ne0,\qquad
 D(0)=153/16.
 \tag{O45P.11}
\]
The nonvanishing assertion is checked by the gcd with \(H\), so it
holds at every one of the twelve marked centres.

The full translation \(T=-G\) on the original source is
\[
 L=\frac{y+\sqrt3/2}{x-R},\qquad
 x_T=L^2-x-R,\qquad y_T=L(x-x_T)-y.
 \tag{O45P.12}
\]
At the image point its plus-\(G\) slope is
\((y_T-\sqrt3/2)/(x_T-R)=-L\). Applying the same chord formula
with this slope returns the original \(x,y\).
On O45P.2 the scaled slope is simply \(t+A\); substituting these
rational source coordinates verifies both composites and the full
source equation without introducing another source family.
At the cusp the denominators are \(c\ne0\). The point is fixed and
the derivatives in the actual graph chart are
\[
 DT=\begin{pmatrix}1&2g_G\\0&1\end{pmatrix},\qquad
 DT^{-1}=\begin{pmatrix}1&-2g_G\\0&1\end{pmatrix},
 \qquad g_G=-\frac{2c}{\sqrt3}.
 \tag{O45P.13}
\]
The derivative of \(t(a,x,y)\) in \(x,y\) is zero at this point, so
these are the derivatives on the full source graph, not on a
different transverse surface. The ordinary addition charts at \(G\)
and \(O\), and the simultaneous infinity charts below, extend the
same mutually inverse translation over the rest of the source.

\par

\medskip\noindent\textbf{O45P.3. The complete point-line coefficient and its two cusp frames}

\par\nobreak

Let \(M_8=\mathcal O([8G]-[O])\), with its canonical rational point
section \(\sigma_8\). The actual tangent/chord construction gives
\[
 f_2=\frac{\ell_G}{x-x(2G)},\qquad
 f_4=f_2^2\frac{\ell_{2G}}{x-x(4G)},\qquad
 f_8=f_4^2\frac{\ell_{4G}}{x-x(8G)},\qquad
 B_{8,1}=f_2^8\,\frac{f_8\circ T}{f_8}.
 \tag{O45P.14}
\]
Every tangent is evaluated at the indicated original section.
The horizontal divisor of \(f_m\) is
\(m[G]-[mG]-(m-1)[O]\). Pullback by \(T\) adds \(G\) to each
section in that divisor. It follows directly that
\[
 \operatorname{div}_{\rm hor}B_{8,1}
      =[G]+[8G]-[9G]-[O]
      =([8G]-[O])-T^*([8G]-[O]).
 \tag{O45P.15}
\]
Thus \(\sigma_8(p)\mapsto B_{8,1}(p)\sigma_8(Tp)\) is the
specified point-module arrow. Its inverse at \(p\) uses
\(B_{8,1}(p+G)^{-1}\); the inverse-source argument is part of the map.

The coefficient through the cuspidal fibre can be computed before
taking a constant. Put \(p=x/y\), so \(x=p^{-2},y=p^{-3}\);
the original \(G\) has additive coordinate \(g_G\), and \(T\) is
\(p\mapsto p-g_G\). The actual tangent at \(jG\) has slope
\(3/(2jg_G)\), and its doubled point is
\(((2jg_G)^{-2},(2jg_G)^{-3})\).
The three chronological tangent quotients give
\[
 f_m(p)=-\frac{m}{g_G^{m-1}}\,
       \frac{(p-g_G)^m}{(p-mg_G)p^{m-1}},
       \qquad m=2,4,8,
 \tag{O45P.16}
\]
and hence
\[
 B_{8,1}(p)=\frac{256}{g_G^8}
       \frac{(p-g_G)(p-8g_G)}{p(p-9g_G)}.
 \tag{O45P.17}
\]
There are two explicitly related frames. On the finite
normalization point/O chart use
\(r(p)=(p-8g_G)/p\) and \(e=\sigma_8/r\).
The arrow on the coefficients of \(e\) is
\[
 B_{8,1}(p)\frac{r(p-g_G)}{r(p)}
          =\frac{256}{g_G^8}.
 \tag{O45P.18}
\]
This equation includes the zero-section and translated section
directions in their actual Cartier modules.

At the cusp \(\sigma_8\) itself is a nonzero frame, since its
horizontal divisor avoids that point. Multiplying the numerator
and denominator of the dimensionless quotient in O45P.17 by
the exact opposite factor \(1+9g_G/p\) gives
\[
 B_{8,1}(x,y)=\frac{256}{g_G^8}
       \frac{1-73g_G^2x+72g_G^3y}{1-81g_G^2x}.
 \tag{O45P.19}
\]
Both numerator and denominator are units at \((x,y)=(0,0)\).
The overlap between this cusp frame and O45P.18 is precisely the
full \(r(p)\). No extension of \(r\) as an invertible cusp-ring
function is used. The reciprocal in the inverse-source chart
extends as well.

The full meromorphic Hom coefficient has no other finite vertical
divisor. Each finite fibre in the disc family is irreducible;
\(T\) preserves it, so \(f_8\circ T\) and \(f_8\) have equal
whole-fibre orders. The original \(G\)-tangent denominator is
the nonzero constant \(\sqrt3\), and its chord quotient \(f_2\)
has whole-fibre order zero. After O45P.15 cancels the actual
horizontal module divisors, the coefficient and its reciprocal
are meromorphic functions without divisorial poles on the smooth
total source. The normal local ring therefore extends both
holomorphically. The explicit nonzero value in O45P.19 also
excludes a parameter-divisor zero at the new II fibre.
This proves the actual total-family point-line map through the
coincident-fibre parameter, with all its module arguments retained.

\par

\medskip\noindent\textbf{O45P.4. The directed infinity Hom degree and every period unit}

\par\nobreak

The original infinity minimal coordinates are
\(z=1/t,\ X=z^2x,\ Y=z^3y\).
The two original source branches at \(G\) have orders three and
five. Indeed \(X_G=1+az^2,\ Y_G=(\sqrt3/2)z^3\), while
\(\sqrt{f_{\min}}=1+az^2+z^3/2+O(z^5)\); the marked branch
cancels the cubic term on exactly one side.
The product has the original discriminant order eight.
Thus \(G\) meets the marked component three of \(I_8\).

Here is an integral verification of the height antecedent used
for the Hom degree. The positive \(A_7\) matrix \(C\) has diagonal
two and adjacent entries minus one. It has determinant eight and
\((C^{-1})_{33}=15/8\). In the ordered actual source classes
\((O,F,C_1,\ldots,C_7,G)\), use
\[
 O^2=G^2=-1,\quad OF=GF=1,\quad GO=0,\quad
 (C_iC_j)=-C,\quad GC_3=1,
 \tag{O45P.20}
\]
with all other unspecified entries zero. The resulting ten-class
Gram determinant is \(-1\). These classes therefore fill the
unimodular rational-surface lattice, and their quotient by the
first nine classes is generated freely by the original \(G\).
Its height is \(2-15/8=1/8\).
For \(jG\) the actual component is \(3j\bmod8\); the height
intersection formula then gives
\[
 G.O=0,\qquad 7G.O=3,\qquad 8G.O=3,\qquad 9G.O=5.
 \tag{O45P.21}
\]
The directed \(-G\) Hom restricted to \(O\) has degree
\(5-0-3-1=1\). Equivalently, the original \(t_G\)-pullback
linearization has degree \(3+1-3+0=1\); these are the two
directions with their own source arguments.

The line \(M_8\) has component degree zero at infinity, since
\(8G\) and \(O\) meet its identity component. The full cycle
intersection matrix has kernel
\(\mathbb Z(1,\ldots,1)\). Its entire remaining vertical
Hom divisor is consequently exactly \(F_\infty\), with
coefficient one fixed by the preceding degree. Multiplication
by the original \(q=t-t_c\) changes it to the selected whole
zero. Thus the original suspension is literally
\[
 B_\lambda(p,v_M)=
       (Tp,\lambda C_*q B_{8,1}(p)v_M).
 \tag{O45P.22}
\]
The continuation in (O45P.22) is defined using the single
nonzero \(C_*\) obtained from the prescribed initial line map
\(L_{a_0}\) at a protected regular point \(s_0\), in the same
Cartier frames:
\[
 C_*=
 \frac{L_{a_0}(s_0)}
 {(t(s_0)-t_c(a_0))B_{8,1,a_0}(s_0)}.
\]
Equality of the divisors makes this ratio an invertible function
on the compact initial source, hence constant on that source.
This proves equality of the entire initial map with
\(C_*(t-t_c)B_{8,1}\). At later marked parameters (O45P.22),
with this same specified \(C_*\), defines the continuation
family whose coefficient maps and inverses extend over the disc.
The boundary identification in this proof concerns this defined
continuation of the initial source, line, cap and reference.

For an independently prescribed varying line map \(L_a\), the
exact comparison retains the analytic evaluation ratio
\[
 c(a)=\frac{L_a(s(a))}
 {(t(s(a))-t_c(a))B_{8,1,a}(s(a))}
\]
wherever these maps and the protected section are defined. Its value
\(c(a_0)=C_*\) alone proves neither constancy nor extension
through a missing parameter. That further family comparison
is not asserted here.

The marked infinity Tate quantities are the full functions
\[
 Q_\infty=z^8\kappa_\infty,\quad
 \kappa_\infty(0)=-a/2304,\qquad
 P_\infty=z^3\pi,\quad \pi(0)=-1/24.
 \tag{O45P.23}
\]
Their constants follow from the inverse \(j\)-coefficient and
the original normalized nodal parametrization
\[
 X=1+\frac{12Z}{(1-Z)^2},\qquad
 \frac Y{\sqrt3}=\frac{12Z(Z+1)}{(Z-1)^3}.
 \tag{O45P.24}
\]
In particular \(b_0=P_\infty^8/Q_\infty^3\) extends and has value
\(-1/(9a^3)\). Substitution in O45P.24 gives the full
identity-component eighth-point values
\[
 X_{8,0}=\frac{81a^6-90a^3+1}{(9a^3+1)^2},\qquad
 \frac{Y_{8,0}}{\sqrt3}
      =\frac{108a^3(9a^3-1)}{(9a^3+1)^3}.
 \tag{O45P.25}
\]
These are also obtained by the three original tangent doublings.
At the II centres \(a^3=3/256\), so \(b_0=-256/27\ne1\).
The point and zero module charts therefore remain separated in
the infinity collar.

For \(f_b(Z)=\Theta_Q(Z/b)/\Theta_Q(Z)\), three actual theta
shifts give
\[
 f_{Q^3b_0}(Z)=(-Z/b_0)^3Q^{-6}f_{b_0}(Z).
 \tag{O45P.26}
\]
The inverse of this entire coefficient changes the original
point coordinate to its identity-component coordinate.
Under \(Z\mapsto P_\infty^{-1}Z\) it changes the line
multiplier by \(P_\infty^3\); under \(Z\mapsto Q_\infty Z\)
it changes \(P_\infty^8\) to \(b_0\).
Thus the full \(Q_\infty^{-6}\) divisor and every base unit
remain in the frame arrow. The valuation columns in this frame
are exactly \((8,0)\) and \((-3,0)\).
Their actual period words and inverses are
\[
 H=h^{-1}B^{-3},\quad K=h^3B^8,\qquad
 h=H^8K^3,\quad B=H^{-3}K^{-1}.
 \tag{O45P.27}
\]
The integer matrix
\(\left(\begin{smallmatrix}-1&3\\-3&8\end{smallmatrix}\right)\)
has determinant one and sends the original source valuations to
\((1,0)\). The second multiplier of \(K\) is \(\lambda^8\)
times its entire extending unit.
The source charts \(X_i=Z/z^i,\ Y_i=z^{i+1}/Z\) have
\(X_iY_i=z\); the complete \(Q_\infty/z^8\) and
\(P_\infty/z^3\) coefficients give their period and translation
overlaps. The identity on the actual point module, combined
with O45P.27, is an isomorphism to the rank-one completed
quotient. Its inverse uses the displayed inverse words and
reciprocal units. These maps include all eight source components.

\par

\medskip\noindent\textbf{O45P.5. The selected collar, prescribed fan and original cap}

\par\nobreak

The selected node stays simple on \(\Delta\). Define
\(\xi(q,u)=\sqrt{f(t_c+q,u)}\) with its marked value in O45P.4,
and put
\[
 h_s(q,u)=\frac{2(g(t_c+q,u)-\xi(q,u)^3)}q,\qquad
 h_s(0,u)=-2u^4.
 \tag{O45P.28}
\]
This is an actual extending unit. The full source chart is
\[
 X=y-(x-\xi)\sqrt{x+2\xi},\qquad
 Y=h_s^{-1}\bigl(y+(x-\xi)\sqrt{x+2\xi}\bigr),\qquad XY=q.
 \tag{O45P.29}
\]
Its inverse sets \(q=XY,\ y=(X+h_sY)/2\) and solves
\(2(x-\xi)\sqrt{x+2\xi}=h_sY-X\). The derivative
\(2\sqrt{3\xi}\) is nonzero, so this is the inverse on the
whole node collar over the source parameter disc.

Keep the full inverse Tate coordinate
\(Q=q\kappa(q,u)\) and \(P=Z(G)\), with their original
reference scales. The actual point-module arrow and periods are
\[
 V_{\rm pt}=f_{P^8}(Z)v_M,\quad
 h:(Z,V_{\rm pt})\mapsto(QZ,P^8V_{\rm pt}),\quad
 B:(Z,V_{\rm pt})\mapsto(P^{-1}Z,\lambda q\beta V_{\rm pt}),
 \tag{O45P.30}
\]
\[
 \beta=C_*B_{8,1}(p)
       \frac{f_{P^8}(P^{-1}Z)}{f_{P^8}(Z)}.
 \tag{O45P.31}
\]
The complete horizontal divisors cancel. On every regular
source fibre this is independent of \(p\); evaluation at a
protected regular point gives the full analytic base function.
The module maps and their reciprocals just proved show that
it is an extending unit. Its value and logarithm are never
replaced by their central limits.

For a prescribed integral basis, each source/line transition
therefore acts on the affine character algebra by its original
monomial times the complete coefficient in O45P.30-\/-31.
For example its two vertex maps, with actual coefficient
characters \(A,D\), are
\[
 (a,b,c)\mapsto((AD)^{-1}a,Ab,Dc),\qquad
 (u,v,w)\mapsto(ADu,D^{-1}v,A^{-1}w).
 \tag{O45P.32}
\]
They preserve the original height product \(q\).
Their reciprocals are the inverse maps, and the same
character evaluation carries every cone of the given fan.
Since \(v\) and all coefficient functions are single-valued
on \(\Delta\), this is the prescribed fan at the centre,
with its unchanged branch order and period arrows.

The original cap disc at \(t=0\) is smooth throughout
\(\Delta\). Continue its complete source and line period
vector \(\gamma(t,u,\lambda)\), with all initial logarithms
and annular integers. Its directed map and inverse are
\[
 \Psi(\zeta,t,u)=
   \left(\zeta+\frac{\log t}{2\pi i}\gamma(t,u,\lambda),t,u\right),
 \qquad
 \Psi^{-1}(\zeta,t,u)=
   \left(\zeta-\frac{\log t}{2\pi i}\gamma(t,u,\lambda),t,u\right).
 \tag{O45P.33}
\]
A logarithm branch change adds precisely the whole original
\(\gamma\)-period in every coordinate. Thus this is the
single-valued map on the original torus quotient.
Its cap meridian maps to the base meridian plus \(\gamma\),
retaining its multiplicity-one nonzero clutch. Its line
coordinate includes the original \(\log\lambda\).
There is no parameter-deck residue on this filled disc:
all the displayed source, point, fan and cap maps return
to themselves by their actual single-valued coefficients.

\par

\medskip\noindent\textbf{O45P.6. Reference sphere, compact family and identity closing}

\par\nobreak

We make the topology antecedent for every prescribed fan
explicit. At the original K10 \(a=b=1\) source, K10\textquotesingle s
integral local-system and cap maps give its standard
six-sphere. The full five source matrices in that proof
have product one; the homology images of \(T_i^{-t}-I\)
generate a primitive rank-three sublattice. Its quotient
is measured by \(\psi\), and the actual cap meridian in
O45P.33 has \(\psi\)-value one. The complete primitive
Gysin/Leray calculation in K10 leaves only \(H^0,H^6\).
This establishes the stated reference sphere from its
actual cap. Its earlier scalar-return value is unused.

For each prescribed lattice-basis fan \(P\), the finite
canonical word of K71/K73 starts at this actual source.
At each adjacent change the common blowup and its two
ruling collapses, with the primitive integral Gysin
and van Kampen maps of F35/K23, preserve simple
connectivity and the homology-sphere groups. The
complete characteristic/period arrows in K71 keep
the same source, line, cap and all full coefficients.
Applying these finitely many actual maps gives the
reference sphere on \(P\); the six-dimensional
homotopy sphere is the standard sphere. This supplies
the fan-specific starting topology before any source
path is taken.

The open ordered cover is connected. Choose a path
from that K10 marking to the chosen boundary point
of \(\Delta\), and cover its compact image by finitely
many small discs avoiding the listed deleted values.
The original source, line, selected height, fixed
fan, infinity charts and cap maps O45P.22-\/-33 apply
on every disc. Their overlaps are the actual source
identities, point modules and integral changes of
period basis. Continue the original cap vector and
all reference integers along the path.

On the compact source outside the selected tube,
choose a smooth Hermitian metric on the actual
point line. The full unscaled lift and its base
source translation have a uniform finite norm over
the compact parameter discs, including the II centre.
Choose a common nonzero scalar whose product with
that bound is less than one. The suspension is
strictly contracting and its inverse expands.
This gives a compact annular fundamental domain.
Only finitely many iterates can meet a fixed compact
subannulus, by the strict norm inequality. Thus the
quotient is proper and Hausdorff, uniformly in the
parameter. A physical-circle equality with an iterate
of \(B\) forces that iterate to be zero by the same
inequality; the quotient line action is the original
physical one.

Inside the selected tube use the actual finite
period representatives of O45P.29-\/-32; at infinity
use the entire cycle and O45P.27; at the ordinary
cap use O45P.33. Their overlaps are the displayed
source/module/period isomorphisms with reciprocal
inverses. The compact collar pieces cover the
fundamental domain and glue to a proper family.
The source projection is a submersion at the cusp
by O45P.10, at the selected fan by its full toric
coordinates, and at the other source/cap pieces
by their continued smooth charts. Therefore the
completed family is a smooth proper holomorphic
submersion over the whole \(\Delta\).

The same construction along the reference path
carries the reference sphere to the disc boundary.
To describe its inverse explicitly, choose local
lifts of the path tangent vector in these charts,
transport them by the displayed differential
overlaps, and patch by a partition of unity on
the compact total space. Their vertical differences
are tangent to the compact fibre. The resulting
horizontal flow exists along the whole path; its
inverse is the reverse flow. It preserves the
orientation induced by the complex fibres.

On the disc itself choose such a connection and
average it over the physical circle. Radial
transport from the centre, with inverse transport
along the reversed radial segment, trivializes
the completed family smoothly over the disc.
On its boundary the analytic closing is exactly
the identity because every source, line, fan,
cap and reference map is single-valued on
\(\Delta\). The boundary return is therefore
isotopic to the identity through these actual
fibre maps. Its physical return class is zero.

Finally take any fixed scalar admitted on the
original boundary family. Decrease its modulus
to the common disc-contraction scalar along the
constant-argument interval. The strict
contraction inequality remains valid on the
boundary throughout this interval. Substitute
the full \(\log\lambda\) in O45P.22,
O45P.27, O45P.30-\/-31 and O45P.33.
This is a proper annular family with the
same endpoint maps and no added scalar turn.
Horizontal flow relative to the endpoint
collars and its reverse give an isotopy of
the two boundary returns. Hence the zero
return holds for every admitted fixed scalar.
This constructs the claimed filling for
each of the twelve marked meridians, with
the original source, cap and prescribed fan.

\par

\medskip\noindent\textbf{O45P.7. Exact computations and geometric scope}

\par\nobreak

The source, Miller and integral-topology computations are
reproduced by
\path{code/current_audits/audit_paper_i/iv_os45/code/k116_afterreading_exact.py}.
The program uses exact rational functions, a full modular
coset enumeration, and Hermite and Smith normal forms;
the 249 checks and all four selected-site matrices are
recorded in
\path{code/current_audits/audit_paper_i/iv_os45/output/k116_afterreading_exact.json}.
The unrestricted fan extension and proper-disc conclusion
follow from the full character, ruling and compactness
maps proved above.

The conclusion concerns the twelve unselected-II source meridians.
The selected cusp, infinity \(I_9\), ordinary-cap collision and
handle/source generators have their separately stated parameter
loops and comparison maps.

% End expanded verified cp36 input



\par\medskip\noindent\textbf{Construction dependencies.}\par
The original K10 source is carried by its literal equation,
point line, selected whole zero and original ordinary cap.
The source and line morphisms above retain every cap and
period coefficient. The compact disc and scalar-interval
flows, with their inverses, give the physical return.

The exact arithmetic uses the displayed inputs and the
program and output named in O45.7 and O45P.7.
The two original Cartier frames, their overlap, the full
infinity theta shift, primitive period inverse, and cap
logarithm occur in the geometric arguments above.


% End expanded verified cp36 input

\end{proof}


\subsection{The complete OS45 selected-cusp return on every fan}
